\begin{abstract}
This thesis investigates the design, implementation, and evaluation of a \ac{LKA} function within a CARLA-based, software-only \ac{VIL} simulation environment. A Python-based control stack is developed to evaluate three progressively richer lateral-control architectures: a cross-track PID, a heading-error PI, and a vision-based PI utilizing a deep convolutional lane detector and \ac{IPM}. The vehicle is modeled as a \ac{DST}, with analytical PID gains derived via system identification and pole placement. This theoretical analysis reveals a structural advantage in the heading-error formulation, which exhibits a gentler gain dependence on vehicle speed ($1/v$) compared to the cross-track approach ($1/v^2$). 

Evaluated on a common test route, the heading-error PI achieves the best performance (\ac{RMSE} $0.120$~m, max error $1.185$~m), outperforming the cross-track PID by nearly a factor of two. The vision-based PI yields an \ac{RMSE} of $0.270$~m—primarily due to steady-state perception biases—but produces smoother physical trajectories over lane discontinuities. Ultimately, this work contributes a reproducible \ac{VIL} implementation, a complete analytical tuning pipeline, explicit junction-handling logic, and a rigorous comparative analysis that lays the groundwork for future extensions like \ac{MPC} and sensor fusion.
\end{abstract}
\acresetall